\section{Limits}\label{sec:limits}

This section collects what bounds the result, with the direction each limit leans where known. It
has four groups: the boundaries of the result, the modelling choices that lean towards or against CPM, the
method, and the inputs and the standing of the evidence.

\subsection{Boundaries of the result}

\emph{The stress boundary.} Under 2 points of distress owners can run on the pool. At 2 points they leave after a
write-off takes their realised return below their alternative, pool capital falls short, and bridges served drop:
on the first development check at $a_{0} = 0.6$ this happened in two of three seeds, with bridges served of 0.29 and
0.30 in years 4 and 5 against 0.90 in the third. The repeated size check counts the seeds with such a run, bridges
served below 0.6 in years 4 and 5, out of eight at each population of 500, 1,000, 2,000, 4,000 and 8,000: none at any
size at $a_{0} = 0.05$; 0, 0, 1, 3 and 3 at $a_{0} = 0.2$; and 1, 4, 7, 7 and 5 at $a_{0} = 0.6$. Where runs occur the counts rise with
population and then level off, and one falls back at the largest size. The run does not go away as the population
grows, so it is a boundary of the design at every size, reached from 20\% launch
participation upward at 2,000 people and more, and populations of 500 to 1,000 hide it. Its cause is long loans funded by capital on two months' notice. Term
lending stays in the model, with the bridge and term parts of the pool's book, income and losses reported
apart. Changing the cash set aside for exits or ranking bridges first after seeing the result would favour the
pool, so neither is done. The limit leans against CPM where it binds. On the size check's three cells at 4,000
people CPM still fails fewer payments than the bank arrangement in all eight seeds at 2 points, and its mean gain
over the bank arrangement and over best-case DeFi is positive in every seed. On the confirmation map the mean gain
is positive in every cell at every level of distress, and CPM fails fewer payments in \CellsFewerFailures{} of the
\CellsAll{} cells (Section~\ref{sec:results}). Figure~\ref{fig:stress} and Table~\ref{tab:stress} show the cells and seeds
with a pool run at each level of distress.

\emph{How a run unfolds.} An owner reviews its holding with a chance of 15.3\% a month, a rate decided for users and
with no source for investors, reads only its own realised return, and does not react to other owners leaving. The
exit notice is fixed at two months. A run where owners leave because others are leaving, the first-mover
advantage of the literature on fund runs, is therefore absent, and the stress boundary is a lower
bound on run risk. This leans toward CPM. Neither the review rate nor the notice is varied.

\emph{Capital and the published quote.} Pool capital holds at its target under the quote rule, because entrants
act on the published quote, which with no fee margin equals their alternative. Under the record rule, when they
read the pool's realised twelve-month return instead, capital sat below 95\% of its target in 29 to 51\% of months on the development checks, at 1,000 and
4,000 people, while the gain over best-case DeFi stayed positive. The quote rule leans toward CPM. The record
rule is run over the confirmation map as the pessimistic one, with its own region. A third rule, each saver
reading the record with the leaving rule's noise, is a side check on three cells. The record rule is one
pessimistic rule, not a proven worst case: savers who trusted the pool less than its own record would do
worse, and on one of the side check's nine numbers the noisy reading serves fewer bridges than it. The two maps
share their seeds and the bank arrangement, so they are one body of evidence read under two
assumptions, not two confirmations.

\emph{The shock, the low rates and the price inputs.} The registered maps hit the pool with one shock, in months
25 to 36, at outside rates of 2\% and above, with an owners' spread of 0.25 point and a custody cost of 0.24\%.
The exploratory side checks of Section~\ref{sec:robustness} move each of these on three seeds and six cells: they
change the sizes; the gains stay positive, and under 2 points of distress the low rates behave as the lowest registered
rate does. A shock in the first months after launch, a rate that moves within a run,
and a rate rise on a book of fixed coupons are not studied.

\subsection{Choices that lean one way}

\emph{Cash advances in the bank arrangement.} Card holders reach the card for bills and wages only through cash
advances. By inference from the model's weights about four fifths of those advances pay wages the holder owes
as an employer, which consumer cash
advances do not pay in practice. On the development check holders' cash advances ran about 15 times the United
States level (\$1,777 a person a year against \$112 to \$118 a consumer unit) and their fees about 31 times (\$163
against \$5.33). The fees avoided were 11 to 27\% of CPM's wealth gain over the bank arrangement and 38 to 50\% of its
gain over best-case DeFi, so the excess leans toward CPM. The card contract is left as built, every gain is
reported with the cash-advance fees and interest saved taken out, and a cell counts in the region only if it passes
both ways.

\emph{The bank arrangement's alternatives.} A payer short under banks and cards can draw on savings, a card
cash advance, a revolved card balance or earned-wage access. Payday loans, buy now pay later and loans from family
are not in the model, and a failed payment is charged the \$27 fee without the bank paying the item, as an
overdraft would. Each gives the bank arrangement fewer ways to cover a shortfall than households have,
which leans toward CPM.

\emph{Best-case DeFi is generous to DeFi.} It keeps the full reserve-funded savings yield, converts to and from bank
money at no fee where $P$ pays about 0.2\%, is accepted by every seller, and bears none of its own
operational or legal risks. Its over-collateralised loan is in the arrangement as it is: a payer who is short has
already used its own money and same-day savings, so nothing is left to lock and the loan covers none of the
shortfall. Each leans against CPM. On the registered clause every group is ahead of best-case DeFi; the
money-only number is the premium paid by people never short, and it decides no verdict.

\emph{No growth.} Total spending is fixed, and no production, growth or wider-economy effect is modelled: a firm that
borrows earns no more afterwards, and completed trade creates no income. This leans against CPM, and is why
completed trade is reported beside the headline as a bracket from zero to face value, not inside it. In the
other direction, a failed bill is erased, not paid later, so completed-trade dollars are an upper bound.

\emph{Owners' return.} Owners' return against their alternative has a loss tail that shrinks with population. One
loan is a share of the pool that falls as one over the number of people: at 1,000 people a single business loan reached 15\% of the
pool and one default drove a shortfall, while at 4,000 owners earned their alternative in every year on an earlier
development check. A loan of \$25,000 to \$100,000 is still 0.7\% of the pool at 8,000. The model also has no permanent default,
since a written-off debt stays owed and is collected later, which keeps bridge losses near zero and leans toward
CPM on the bridge price. The loan-loss bound prices the published charge-off rates instead. Owners' return against
their alternative is reported with its interval and read by no verdict; a separate falsifier fires only if owners'
return is shown below DeFi's savings rate. On the confirmation map owners' return less their alternative has a
median within \OwnersMedianWorst{} of a percentage point of zero at every outside rate, and one seed-cell in ten falls
short by \OwnersTenthLeast{} to \OwnersTenthWorst{} of a point or more (Table~\ref{tab:owners}).

\subsection{Method}

\emph{A paired transient.} Each run lasts 60 months after launch and every gain is measured during a transition:
participation is still growing in most cells when the run ends, so the comparison is of two paths from the same
start over the same five years, not of two steady states. The lowest launch level is 5\%, so neither a launch from
an empty network nor a long-run state is shown. The direction of this limit is not known. The mean wealth gain is
also read at 36 and 48 months (Table~\ref{tab:gains}), and Figure~\ref{fig:participation} shows participation still
rising at month 60 from every launch level.

\emph{Settlement within a day.} The model's day has no earlier and later. Payments due on a day are settled
together, money received can be paid on within the day, and where more than one outcome fits the
day's rules the one with the most payments completed is taken (Proposition~\ref{prop:settlement}). This is the same
in every arrangement and leans toward completion in each. The payments that rest on the choice of outcome are at
most \NettingShareMost{} of completed value in any run, so the comparison between arrangements does not turn on
it. The daily clock removes failure through timing within a day, from every arrangement alike.

\emph{A day the settlement cannot resolve.} When the pool's cash is short its requests are served pro rata and the
day's map need not be monotone (Section~\ref{sec:updating}). The model resolves two states that repeat when a
request qualifies in one and fails the repayment check in the other. A day when the iteration still does not
come to rest has no further rule: the model stops the run and does not guess. This happened in no run of the two
confirmation maps or the side checks. It happened in \CutSeedsAside{} of the \CutSeedsTried{} seeds of
the rate cut, once each, at outside rates of 3\% and 3.5\%; each such seed is set aside whole and named in the
cut's record. The cut is therefore read on the seeds where every run completed, and what those four seeds would
have shown is not known.

\emph{Chains of failure.} Failure passed on within a day is measured exactly. Failure passed on across days is not:
the spread of a shock to people not laid off mixes the payment network with the money that extra layoffs'
wages take out of the household economy and with the shift in everyone's spending. A count of chains across days
would need a record of receipts lost by payee and a run where those wages stay inside the economy.

\emph{What demurrage adds on its own.} CPM against the pool without the charge, on the same people and draws, is
what the holding charge adds. This edge has not settled with population size: it grows at every doubling from 500 to
8,000 people. With completed purchases at their price, as the size check reads it, between 4,000 and 8,000
it moves from \$15.4 to \$19.0 a person a year at $a_{0} = 0.05$, from
\$35.6 to \$43.5 at $a_{0} = 0.2$ (inside its interval) and from \$64.8 to \$78.8 at $a_{0} = 0.6$. At 4,000 people the paper therefore understates what the
charge adds, which leans against CPM.

\emph{The adoption rule and the network.} The base joins people by a simple contagion among trading partners drawn
uniformly, with a scattered launch. Section~\ref{sec:robustness} replaces each on the development seeds, and the sign
of CPM's gain survives every replacement while the speed of adoption does not. Thresholds for a payment money have
not been measured: the measured ones are those of one paid online service, and a common threshold is a design, not
an estimate. Under a common threshold of 0.3 a launch of 10\% or less does not tip within five years, and the gain
per person is then the smaller gain of a network that stays small. The hub world's tail is that of the number of
trading partners in one community currency, whose source could not estimate the tail of amounts; a heavier tail, a
network grown from payment data and one that changes as people join are not modelled. Wage links are uniform in
every world.

\subsection{Inputs and evidence}

\emph{Inputs and their sources.} The model's inputs are United States figures, each listed with its value and
source in Appendix~\ref{app:parameters}, as in the agent-based models used in central banks. The question is
whether CPM moves money better in that economy; no output forecasts a quantity of the United States economy. Some
inputs are proxies or have no primary source: the 1.7\% merchant card fee, found only in trade
press, which scales the saving against the bank arrangement and leaves the comparison with best-case DeFi unchanged;
the issuer's cost of 0.24\% of $P$, a money-market fund expense ratio; the 43\% debt cap; and the loan terms.
Mortgages are left out. No unemployment insurance is modelled, which overstates the losses of workers hit by the
stress and leans against the pool, since benefits would be receipts a bridge lends against.

\emph{Inputs set by declaration.} The parameter table marks each input as sourced or declared. Of those
declared or resting on a proxy, the leaving noise, the fixed cost, the failure-fee cap, the buffer review share, the
custody cost, the owners' spread and the date of the shock have a sensitivity run. The shape and median of initial wealth, the
spending elasticity, the wage share, the number of purchases a month, the twelve-month windows, the debt-to-income
limit, the exit notice and the review rate have none. The direction in which each would move the result is not
known.

\emph{Where the robustness checks sit.} The checks of Section~\ref{sec:robustness} other than the two bound maps and
the spread of distress are run at the central outside rate of 4.4\%: the side checks, the finer cut of the stress
boundary and the adoption, network and launch variants on development seeds, the exploratory side checks on three
side seeds, and the size check on its own eight. The two bound maps cover all three rates. Nothing here shows that the variants
leave the result unchanged at 2\% or 6\%.

\emph{What registered means.} The claim, its tests and the seed count were fixed in the dated record of decisions
and the analysis plan shipped with the code, before any confirmation seed was run. They were not deposited with a
public registry, so their dates rest on that record. The plan lists its departures from earlier registered
text, and the two additions made after the first confirmation map was read, the rate cut and the exploratory side
checks, are labelled exploratory there and here.

\emph{Development and confirmation.} The headline results and the registered claim are read on the confirmation
map, \SeedsConfirm{} seeds never run before the freeze and read once after it. The loan-loss bound, the
interest-returned bound, the side checks, the finer cut of the stress boundary, the robustness sets, the
update-scheme reading and the spread of distress are read on the three development seeds, on which the model was
built and its rules decided, and the size check on eight seeds of its own. Each is labelled so, and none is
confirmed.
